\documentclass[11pt]{aghdpl}
% \documentclass[en,11pt]{aghdpl}  % praca w języku angielskim

% Lista wszystkich języków stanowiących języki pozycji bibliograficznych użytych w pracy.
% (Zgodnie z zasadami tworzenia bibliografii każda pozycja powinna zostać utworzona zgodnie z zasadami języka, w którym dana publikacja została napisana.)
\usepackage[english,polish]{babel}

% Użyj polskiego łamania wyrazów (zamiast domyślnego angielskiego).
\usepackage{polski}

\usepackage[utf8]{inputenc}

% dodatkowe pakiety

\usepackage{mathtools}
\usepackage{amsfonts}
\usepackage{amsmath}
\usepackage{amsthm}

% --- < bibliografia > ---

\usepackage[
style=numeric,
sorting=none,
%
% Zastosuj styl wpisu bibliograficznego właściwy językowi publikacji.
language=autobib,
autolang=other,
% Zapisuj datę dostępu do strony WWW w formacie RRRR-MM-DD.
urldate=iso8601,
% Nie dodawaj numerów stron, na których występuje cytowanie.
backref=false,
% Podawaj ISBN.
isbn=true,
% Nie podawaj URL-i, o ile nie jest to konieczne.
url=false,
%
% Ustawienia związane z polskimi normami dla bibliografii.
maxbibnames=3,
% Jeżeli używamy BibTeXa:
backend=bibtex
]{biblatex}

\usepackage{csquotes}
% Ponieważ `csquotes` nie posiada polskiego stylu, można skorzystać z mocno zbliżonego stylu chorwackiego.
\DeclareQuoteAlias{croatian}{polish}

\addbibresource{bibliografia.bib}

% Nie wyświetlaj wybranych pól.
%\AtEveryBibitem{\clearfield{note}}


% ------------------------
% --- < listingi > ---

% Użyj czcionki kroju Courier.
\usepackage{courier}

\usepackage{listings}
\lstloadlanguages{TeX}

\lstset{
	literate={ą}{{\k{a}}}1
           {ć}{{\'c}}1
           {ę}{{\k{e}}}1
           {ó}{{\'o}}1
           {ń}{{\'n}}1
           {ł}{{\l{}}}1
           {ś}{{\'s}}1
           {ź}{{\'z}}1
           {ż}{{\.z}}1
           {Ą}{{\k{A}}}1
           {Ć}{{\'C}}1
           {Ę}{{\k{E}}}1
           {Ó}{{\'O}}1
           {Ń}{{\'N}}1
           {Ł}{{\L{}}}1
           {Ś}{{\'S}}1
           {Ź}{{\'Z}}1
           {Ż}{{\.Z}}1,
	basicstyle=\footnotesize\ttfamily,
}

% ------------------------

\AtBeginDocument{
	\renewcommand{\tablename}{Tabela}
	\renewcommand{\figurename}{Rys.}
}

% ------------------------
% --- < tabele > ---

\usepackage{array}
\usepackage{tabularx}
\usepackage{multirow}
\usepackage{booktabs}
\usepackage{makecell}
\usepackage[flushleft]{threeparttable}

% defines the X column to use m (\parbox[c]) instead of p (`parbox[t]`)
\newcolumntype{C}[1]{>{\hsize=#1\hsize\centering\arraybackslash}X}


%---------------------------------------------------------------------------

\author{Marcin Szpyrka}
\shortauthor{M. Szpyrka}

%\titlePL{Przygotowanie bardzo długiej i pasjonującej pracy dyplomowej w~systemie~\LaTeX}
%\titleEN{Preparation of a very long and fascinating bachelor or master thesis in \LaTeX}

\titlePL{Przygotowanie pracy dyplomowej w~systemie~\LaTeX}
\titleEN{Thesis in \LaTeX}


\shorttitlePL{Przygotowanie pracy dyplomowej w~systemie \LaTeX} % skrócona wersja tytułu jeśli jest bardzo długi
\shorttitleEN{Preparation of a long and fascinating thesis in \LaTeX}

\thesistype{Praca dyplomowa magisterska}
%\thesistype{Master of Science Thesis}

\supervisor{prof. dr hab. Marcin Szpyrka}
%\supervisor{Marcin Szpyrka PhD, DSc}

\degreeprogramme{Informatyka}
%\degreeprogramme{Computer Science}

\date{2015}

\department{Katedra Informatyki Stosowanej}
%\department{Department of Applied Computer Science}

\faculty{Wydział Elektrotechniki, Automatyki,\protect\\[-1mm] Informatyki i Inżynierii Biomedycznej}
%\faculty{Faculty of Electrical Engineering, Automatics, Computer Science and Biomedical Engineering}

\acknowledgements{Serdecznie dziękuję \dots tu ciąg dalszych podziękowań np. dla promotora, żony, sąsiada itp.}


\setlength{\cftsecnumwidth}{10mm}

%---------------------------------------------------------------------------
\setcounter{secnumdepth}{4}
\brokenpenalty=10000\relax

\begin{document}

\titlepages

% Ponowne zdefiniowanie stylu `plain`, aby usunąć numer strony z pierwszej strony spisu treści i poszczególnych rozdziałów.
\fancypagestyle{plain}
{
	% Usuń nagłówek i stopkę
	\fancyhf{}
	% Usuń linie.
	\renewcommand{\headrulewidth}{0pt}
	\renewcommand{\footrulewidth}{0pt}
}

\setcounter{tocdepth}{2}
\tableofcontents
\clearpage

\chapter{Inconsistency indexes}
\label{sec:inconsistencyIndexes}
\section{Inconsistency indexes for complete matrices}

This subsection presents sixteen common inconsistency indexes. Their detailed description, including the formulas, is necessary to modify them in the next step so that they can also work for incomplete matrices. Many of them were described and tested numerically in article \cite{Brunelli2013}. In all methods, it is assumed that the PC matrix is reciprocal.


 \subsection{Saaty index ($CI$)}

This is one of the most important and popular indexes and was introduced by Saaty \cite{SAATY1977}. In order to determine inconsistency, the matrix's eigenvalues should be computed. The author used the dependence that the largest eigenvalue of the matrix is equal to its dimension if and only if the given matrix is completely consistent. On this assumption, he based his thoughts and proposed the formula:
	\begin{equation} 
		CI(M)=\frac{\lambda_{max}-n}{n-1},
	 \end{equation}
 where $\lambda_{max}$
  is the principal eigenvalue of the PC matrix and $n$
  is its dimension.


\subsection{Geometric consistency index ($GCI$)}

This index on the assumption ~\ref{eq:consistent2} was proposed by Craford and Williams \cite{CRAWFORD1985} and then refined by Aguaròn and Moreno-Jimènez \cite{AGUARON2003}. In this case the priority vector should be calculated using the geometric mean method. Consider~\ref{eq:consistent} one can create a matrix:
	\begin{equation} 
		E=\left[e_{ij}\mid e_{ij}=m_{ij}\frac{w_{j}}{w_{i}}\right],\,\,\,\,\,\,i,j=1,...,n.
	 \end{equation}
 The inconsistency index is calculated as follows:
	 \begin{equation} 
		\label{eq:GCI}
		GCI=\frac{2}{(n-1)(n-2)}\sum_{i=1}^{n}\sum_{j=i+1}^{n}ln^{2}e_{ij}.
	 \end{equation}
 

\subsection{Koczkodaj index ($K$)}

One of the most popular inconsistency indexes was proposed by Koczkodaj \cite{KOCZKODAJ1993}. It is based directly on the definition of consistency~\ref{eq:consistent}. The value of the inconsistency index for one triad triad was defined as:
	\begin{equation} 
		\label{eq:K}
		K_{i,j,k}=min\{\frac{1}{m_{ij}}\mid m_{ij}-\frac{m_{ik}}{m_{jk}}\mid,\frac{1}{m_{ij}}\mid m_{ik}-m_{ij}m_{jk}\mid,\frac{1}{m_{jk}}\mid m_{jk}-\frac{m_{ik}}{m_{ij}}\mid\}.
	 \end{equation}

 This formula was simplified by Duszak and Koczkodaj \cite{DUSZAK1994} and is given as:
 	\begin{equation} 
		K(\alpha,\beta,\gamma)=min\{\mid1-\frac{\beta}{\alpha\gamma}\mid,\mid1-\frac{\alpha\gamma}{\beta}\mid\},\,\,\,\,\,\,where\,\alpha=m_{ij},\beta=m_{ik},\gamma=m_{jk}
	 \end{equation}
 Then it was generalized \cite{DUSZAK1994} for $n>2$. Finally, the inconsistency index has the following form:
 	\begin{equation} 
		K=max\{K(\alpha,\beta,\gamma)|1\leq i<j<k\leq n\}.
	 \end{equation}
 It is worth noting that not only does the coefficient find the greatest inconsistency but also indicates the place in which it occurs.


\subsection{Kazibudzki indexes ($MLTI, MLTI^{*}, CMLTI*$)}

Based on the Koczkodaj inconsistency index and observation that $ln(\frac{\alpha\gamma}{\beta})=-ln(\frac{\beta}{\alpha\gamma})$, Kazibudzki proposed several additional inconsistency indexes \cite{Kazibudzki2016}. Instead of the formula for inconsistency of the triad [eq:k-abg], he introduced two new formulas:
	\begin{equation} 
		LTI(\alpha,\beta\gamma)=\mid ln(\frac{\alpha\gamma}{\beta})\mid,
	 \end{equation}
	\begin{equation}
		\label{eq:lti*} 
		LTI*(\alpha,\beta\gamma)=ln^{2}(\frac{\alpha\gamma}{\beta}).
	 \end{equation}
Based on the above equations, Kazibudzki proposed new indexes. The simplest ones use the geometric mean of the triads. Thus, new indexes could be written in the form:
	\begin{equation} 
		MLTI(LTI)=\frac{1}{n}\sum_{i=1}^{n}\left[LTI_{i}(\alpha,\beta\gamma)\right],
	 \end{equation}
 	\begin{equation} 
		MLTI(LTI*)=\frac{1}{n}\sum_{i=1}^{n}\left[LTI*_{i}(\alpha,\beta\gamma)\right].
			 \end{equation}
 

After further research Kazibudzki introduces another inconsistency index \cite{Kazibudzki2017}, again based on~\ref{eq:lti*}. It was defined as $CM(LTI*)=\frac{MEAN[LTI*(\alpha,\beta,\gamma)]}{1+MAX[LTI*(\alpha,\beta,\gamma)]}$. Hence,
	\begin{equation} 
		CM(LTI*)=\frac{\frac{1}{n}\sum_{i=1}^{n}[LTI*_{i}(\alpha,\beta,\gamma)]}{1+max\{LTI*_{i}(\alpha,\beta,\gamma)\}}.
	 \end{equation}
 

\subsection{Index of determinants ($PL$)}

This index was proposed by Pelaez and Lamata \cite{PELAEZ2003} and is also based on the concept of triad. The authors noticed that PC  matrices can be constructed on the basis of triads. Their determinant is closely related to the consistency of the matrix.

For every triad $(m_{ik},m_{ij},m_{jk})$ one can build a matrix in the form: 
	\begin{equation} 
		T_{ijk}=\left(\begin{array}{ccc}
			1 & m_{ij} & m_{ik}\\
			\frac{1}{m_{ij}} & 1 & m_{jk}\\
			\frac{1}{m_{ik}} & \frac{1}{m_{jk}} & 1
		\end{array}\right),\,\,\,\,\,\,where\,i<j<k.
	\end{equation}
 The determinant of this matrix is:
	\begin{equation} 
		det(M)=\frac{m_{ik}}{m_{ij}m_{jk}}+\frac{m_{ij}m_{jk}}{m_{ik}}-2.
	 \end{equation}
 If the matrix is fully consistent, then $det(M)=0$, else $det(M)>0$. Based on the above considerations, the authors introduced the new inconsistency index that can be formulated as follows:
 	\begin{equation} 
		CI*=\frac{1}{n}\sum_{i=1}^{n}\left(\frac{m_{ik}}{m_{ij}m_{jk}}+\frac{m_{ij}m_{jk}}{m_{ik}}-2\right).
	 \end{equation}
 

\subsection{Kułakowski and Szybowski indexes ($I_1, I_2, I_{\alpha}, I_{\alpha,\beta}$ ) }

Kułakowski and Szybowski proposed two further inconsistency indexes \cite{KULAKOWSKI20141}, which are also based on triads. They use the fact that the number of triads that can be found in a PC matrix is:
	\begin{equation} 
		\label{eq:nPo3}
		\binom{n}{3}=\frac{n!}{(n-3)!3!}=\frac{n(n-1)(n-2)}{6}.
	 \end{equation}
 The index is formulated as follows:
	 \begin{equation} 
	 \label{eq:I1}		
		I_{1}=\frac{6\sum_{t\in T}K(t)}{n(n-1)(n-2)},
	 \end{equation}
 where $K(t)$ is the Koczkodaj index for triad $t=(\alpha,\beta,\gamma)$ of the set of all triads $T$. 

The second inconsistency index is similar:
	\begin{equation} 
	 \label{eq:I2}				
		I_{2}=\frac{6\sqrt{\sum_{t\in T}K^{2}(t)}}{n(n-1)(n-2)}.
	 \end{equation}

Indexes can be combined with each other to create new coefficients. In this way Kułakowski and Szybowski proposed two new indexes. The first one is based on~\ref{eq:K} and~\ref{eq:I1}. This index allows to choose which value should have more impact on the result: the greatest inconsistency found or the average inconsistency of all triads. The new inconsistency index looks as follows:
	\begin{equation} 
		I_{\alpha}=\alpha K+(1-\alpha)I_{1},
	 \end{equation}
 where $0\leq\alpha\leq1$.
  
The second index expands the first one by~\ref{eq:I2}:
	\begin{equation} 
		I_{\alpha,\beta}=\alpha K+\beta I_{1}+(1-\alpha-\beta)I_{2}.
	 \end{equation}
 

\subsection{Harmonic consistency index ($HCI$)}

This index was introduced by Stein and Mizzi and it presents a completely new method of inconsistency computing \cite{STEIN2007}. At the beginning it requires the creation of an auxiliary vector $s=(s_{1},...,s_{n})^{T}$, where $n$ is the dimension of the matrix $M$, for which the index will be calculated. Each element of the vector $s$ is the sum of values in one column of the matrix $M$. Hence, 
	\begin{equation} 
		s_{j}=\sum_{i=1}^{n}m_{ji}\,\,\,\,\,\,\forall j.
	 \end{equation}
 The authors proved that if the matrix $M$ is consistent, then $\sum_{j=1}^{n}s_{j}^{-1}=1$. The formula for the mean harmonic looks as follows (Brunelli, 2015):
	 \begin{equation} 
		\label{eq:hm1}
		HM=\frac{n}{\sum_{j=1}^{n}\frac{1}{s_{j}}}.
	 \end{equation}
 The final formula for the inconsistency index was obtained by normalizing the above equation~\ref{eq:hm1}:
 	\begin{equation} 
		HCI=\frac{\left(HM(s)-n\right)\left(n+1\right)}{n(n-1)}.
	 \end{equation}
 

\subsection{Golden and Wang index ($GW$)}

This index was introduced by Golden and Wang \cite{Golden1989}. It assumes that the priority vector was calculated using the geometric mean method, then normalized to add up to $1$. In this way vector $g*=[g{}_{1,}^{*},...,g_{n}^{*}]$ was obtained, where $n$ is the dimension of the matrix $M$. The next step is to normalize each column of the matrix $M$. After this, the sum of the elements of each column in matrix $M$ is $1$. The obtained matrix is marked with the symbol $M^{*}$. The inconsistency index is defined as follows:
	\begin{equation} 
		GW=\frac{1}{n}\sum_{i=1}^{n}\sum_{j=1}^{n}\mid m_{ij}^{*}-g_{i}^{*}\mid.
	 \end{equation}
 

\subsection{Salo and Hamalainen index ($CM$)}

The index proposed by Salo and Hamalainen \cite{SALO1995} uses the definition of inconsistency~\ref{eq:consistent}. However, it requires the creation of an auxiliary matrix, in which each element contains the smallest and largest discrepancy from consistency based on the formula~\ref{eq:consistent}. The index takes all triads into account:
	\begin{equation} 
		R=(r_{ij})_{nxn}=\left(\begin{array}{ccc}
			[\underline{r}_{11},\overline{r}_{11}] & \ldots & [\underline{r}_{1n},\overline{r}_{1n}]\\
			\vdots & \ddots & \vdots\\{}
			[\underline{r}_{n1},\overline{r}_{n1}] & \ldots & [\underline{r}_{nn},\overline{r}_{nn}]
		\end{array}\right),
	\end{equation}
 where $\underline{r_{ij}}=min\left\{ m_{ik}m_{kj}\mid k=1,\ldots,n\right\}$ , $\overline{r_{ij}}=max\left\{ m_{ik}m_{kj}\mid k=1,\ldots,n\right\}$ and $n$ is the dimension of the tested matrix $M$. A numerical example was presented in \cite{Brunelli2015}. Based on the resulting matrix $R$, the authors proposed the following inconsistency index:
 	\begin{equation} 
		CM=\frac{2}{n(n-1)}\sum_{i=1}^{n-1}\sum_{j=i+1}^{n}\frac{\overline{r}_{ij}-\underline{r}_{ij}}{\left(1+\overline{r}_{ij}\right)\left(1+\underline{r}_{ij}\right)}.
	 \end{equation}
 

\subsection{Cavallo and D’Apuzzo index ($I_{CD}$)}

The authors Cavallo and D'Apuzzo based their index on triads but they conducted studies on a new path, generalizing them for linear, ordered abelian groups (\cite{Cavallo2009}, \cite{Cavallo2010}). Thanks to this, the index can be used also with other relations \cite{Brunelli2013}. Index for relation $max$ can be presented in the form of a formula:
	\begin{equation} 
		\label{eq:CavDAp}
		I_{CD}=\prod_{i=1}^{n-2}\prod_{j=i+1}^{n-2}\prod_{k=j+1}^{n}\left(max\left\{ \frac{m_{ik}}{m_{ij}m_{jk}},\frac{m_{ij}m_{jk}}{m_{ik}}\right\} \right){}^{\frac{1}{\binom{n}{3}}}.
	 \end{equation}
 

\subsection{Relative error ($RE$)}

This index, proposed by Barzaili \cite{Jonathan1998}, requires calculation of the weight vector using the arithmetic mean method for each row and creation of two additional matrices. Thus, the weight vector is $$w_{i}=\frac{1}{n}\sum_{j=1}^{n}m_{ij},$$ where $n$ is the dimension of the matrix. The two auxiliary matrices are calculated according to the formulas:
$$C=\left(c_{ij}\right)=\left(w_{i}-w_{j}\right)$$
$$E=\left(e_{ij}\right)=\left(m_{ij}-c_{ij}\right)$$

Ultimately, the formula for the Relative error is the following:
	\begin{equation} 
		RE(M)=\frac{\sum_{ij}e_{ij}^{2}}{\sum_{ij}m_{ij}^{2}}.
	 \end{equation}


\section{Inconsistency indexes for incomplete matrices}
\label{sec:inconsistencyIndexesForIncompleteMatrices}

There are no inconsistency indexes for incomplete matrices. However, Those presented in chapter 3 could be used in such cases. It requires usually a slight modification of the index definition or calculation only for selected data. The ways in which the examined indexes were adjusted to be able to deal with incomplete matrices are presented below.

\begin{description}

\item[Saaty index] \hfill \\ 
	The input matrix is modified using the method proposed by Harker \cite{HARKER1987}. It means that values $c+1$, where $c$ is the number of non-empty elements in a given row, are placed on the diagonal.

\item[Geometric consistency index]: \hfill \\
	During calculating the weight vector by the geometric mean, empty values are omitted. Additionally, in the formula~\ref{eq:GCI} only non-empty elements $e_{ij}$ are used. The reason for this exclusion is that the domain of the logarithmic function is $R^{+}$.

\item[Koczkodaj index, Kazibudzki indexes, Index of determinants:] \hfill \\ 
  Only those triads which do not contain empty values are taken into account.

\item[Kułakowski and Szybowski indexes]: \hfill \\ 
	Only those triads which do not contain empty values are taken into account. In addition, the number of triads is no longer calculated according to the formula~\ref{eq:nPo3} but determined directly by counting the number of triads.

\item[Harmonic consistency index:] \hfill \\ 
  No modification.

\item[Golden and Wang index:] \hfill \\ 
  During calculating the weight vector by the geometric mean empty values are omitted.

\item[Salo and Hamalainen:] \hfill \\ 
  No modification.

\item[Cavallo and D'Appuzo:] \hfill \\
	During calculating the product~\ref{eq:CavDAp} empty values are omitted.

\item[Relative index:] \hfill \\ 
  No modification.
\end{description}

Proposed methods of adjustments of indexes allow to apply them in incomplete matrices. However, they do not guarantee the best results. It means that further experiments are required. The existing indexes can be modified in many different ways. In this paper emphasis is put on the fact that the adjustments were slight. The aim of this study is to examine existing indexes, not to create new ones.
\chapter{Introduction}
\label{cha:wprowadzenie}

\section{Pairwise Comparisons method}
\label{sec:metodaPorowan}
People have made decisions for ages. Some of them are very simple and come easily but others, more complicated, require deeper analysis. It happens when there are many compared objects, which are complex and the selection criterion is hard to measure precisely. Fortunately, the development of mathematics brought an interesting tool - \textit{The Pairwise Comparisons (PC) Method}. The first case of using the method (in a very simple version) is the election system described by Ramond Llull \cite{Colomer2013} in the thirteenth century. Its rules were based on the fact that the candidates were pairwise compared with each other and the winner was the one who won in the largest number of direct comparisons. The method was reinvented in the eighteenth century by Condorcet and Bord \cite{Kulakowski2016} as they proposed it in their voting system. In the twentieth century, the method found the application in the theory of social choice, the main representatives of which were the Nobel prize winners Keneth Arrow \cite{Arrow} and Amartya Sen \cite{Sen}. The current shape of the method was influenced by the changes introduced by Fechner \cite{Fechner1966} and then refined by {Thrustone} \cite{Thurstone1994}. However, the breakthrough was the introduction to the method \textit\textit{The Analytic Hierarchy Process (AHP)} by Saaty \cite{Saaty2008}, which allowed to compare many more complex objects and create a hierarchical structure.The primary aim of this paper is to check which method of calculating the inconsistency is the best in this case. In order to do it, a series of tests were carried out on various known inconsistency indexes, taking into account many different parameters: the matrix size, the amount of missing data and the level of inconsistency. The results of the research are included in this paper.

The PC method is based on the assumption that it is not worth comparing all objects at the same time. It is better to compare them in pairs and then gather the results together. Such pairwise comparisons are much more intuitive and natural for a human being. How can one be sure that these judgments are consistent? Or what to do if some comparisons are missing? In such a case, is it worth taking the PC method at all?

The answer to the first question is the concept of inconsistency introduced into the method. This paper tries to answer the next two questions - meaning to examine whether available methods for determining inconsistencies give reliable results when a part of the comparisons is missing. 

\section{The aim of the work}
\label{sec:celePracy}
The main aim of this paper is to check which method of calculating the inconsistency is the best when some comparisons are missing. In order to do so, a series of tests was carried out on various known inconsistency indexes, taking into account many different parameters: the matrix size, the amount of missing data and the level of inconsistency.
The tests were implemented in the \textit{R} language which properly fit into numerical calculations.

It was not known from the beginning what the result of this work would be. It was considered that one or more existing inconsistency indexes would turn out to be appropriate for incomplete matrices also or that the tests would show that the inconsistency could not be calculated by these indexes. 
The results of the research are included in this study.

\section{Content of the work}
\label{sec:zawartoscPracy}
The work includes the theoretical part and the description of conducted experiments and their results. It consists of six chapters.\\
The second chapter shows the Pairwise Comparisons method and the problem of inconsistency. Understanding the basics is necessary to go into the further chapters.\\
In the third chapter, sixteen available methods for calculating the inconsistency for the PC matrices are presented. \\
The fourth chapter shows ideas which allow to perform tests. These are modifications of existing inconsistency indexes which enable to adjust them to incomplete matrices and the algorithm which tests the quality of the modified indexes. \\
The results of experiments are presented and discussed in the fifth chapter.\\
The last chapter contains summary, conclusions and ideas for future researches.
\chapter{Pairwise Comparisons method}
\label{sec:pcMethod}
  \section{PC matrix}
	\label{subsec:macierzPC}
	
	The \textit{PC method} is used to choose the best alternative from a set of concepts. However, this goal is achieved by comparing in pairs. A numerical value is assigned to each pair. It not only determines which alternative is preferred but also informs about the strength of this preference. In this way, the finite set of concepts $C=\left\{ c_{1},\ldots,c_{n}\right\} $ is transformed into a \textit{PC matrix} $M=\left(m_{ij}\right)$, where $m_{i,j}\in \mathbb{R}$ and $i,j\in\left\{ 1,\ldots,n\right\}$. The PC matrix for $n$ concepts is following:
$$
M = 
\left(
\begin{array}{lllll}
	1 & m_{12} & \dots & m_{1n}\\
	m_{21} & 1 & \dots & m_{2n}\\
	\vdots & \vdots & \ddots & \vdots\\
	m_{n1} & m_{n2} & \dots & 1\\ 	
\end{array}
\right)
$$

	It is worth noting that the values $m_{ij}$ and $m_{ji}$ represent the same pair. Therefore, one should expect that $m_{ji}=\frac{1}{m_{ij}}$. If
	\begin{equation} 
		\forall _{i,j\in\left\{ 1,\ldots,n\right\}} :m_{ij}=\frac{1}{m_{ji}},
	\end{equation}
		the matrix is called a \textit{reciprocal}.

  \section{Weight vector}
	\label{subsec:wektorWag}
	
	The PC matrix is the basis for calculating the method. It is used in a function $\mu:C\rightarrow \mathbb{R}$ that assigns a positive real number to each alternative in the set $C$. The vector $$\mu=\left[\mu\left(c_{1}\right),\ldots,\mu\left(c_{n}\right)\right]$$ formed in this way is called \textit{the weight vector} or \textit{the priority vector} (see Fig.2.1). It informs which alternative won.

\begin{figure}[ht]
\centerline{\includegraphics[scale=2.5]{fig1.png}}
\caption{Diagram of calculating \textit{the weight vector} based on the article \cite{Kulakowski2014}}
\end{figure}

There are many ways to calculate the vector $\mu$. Among the popular ones are the method using the matrix's eigenvalues or the method based on geometric means \cite{SAATY1998}.

\section{Inconsistency}
\label{subsec:inconsistency}
The second important parameter describing the PC matrix is \textit{consistency}. The matrix is consistent if 
	\begin{equation}
		\label{eq:consistent}
		\forall _{i,j,k \in\left\{ 1,\ldots,n\right\}} :m_{ik}=m_{ij}m_{jk}.
	\end{equation}
Three numbers that should meet this assumption are called a \textit{triad}.

If the matrix is consistent and the weight vector $\mu$ is computed, then for each of the variables $i,j$, where $1\leq i,j\leq n$ meets the equation: 
	\begin{equation} 
		\label{eq:consistent2}		
		m_{ij}=\frac{\mu_{i}}{\mu_{j}}.
 	\end{equation}
In practice, it is very rare for a matrix $M$ to be completely consistent. In the long history of the PC method, many methods were developed to calculate inconsistencies. Many of them are based directly on the definition of consistency~\ref{eq:consistent}, some methods use the eigenvalues of the matrix, others are based on the assumption that each fully consistent matrix fulfills the condition~\ref{eq:consistent}.
\chapter{Testy}

\section{Test URL-a}

Wejdź na stronę \url{https://www.google.pl/} i wpisz szukane zdanie.

\clearpage

\section{Test dzielenia wdów}

Lorem ipsum dolor sit amet, ex est alia dolorem commune. Duo modo errem no. Ea harum doming atomorum mei. Consul animal malorum cu qui, sumo dicta graece an est, vim ei clita regione.

Vel eu quando doming fastidii, mei graeco indoctum an, legere theophrastus in pro. Te mei probatus eleifend interpretaris. Est no autem liber vituperatoribus, cu mea dicam constituto. Ea laudem tritani consectetuer sit, sanctus patrioque expetendis vix in. Duo id fugit adversarium signiferumque, an quot modus molestiae qui.

Ut paulo definiebas pro. Mea an quod esse. Et atomorum facilisis moderatius sit. Graeco iudicabit forensibus in vel. Eam cu lorem aeterno offendit, cu vix nulla congue posidonium. Vel lucilius evertitur vituperata no.

Mea eu graecis prodesset. Et tota eius nec. Ei etiam oratio has, vel ei homero eripuit invenire. Sed ex errem intellegebat, sea et elitr intellegat constituto. Nostro voluptua accusamus eos in, ei sale admodum has. Vim ne consetetur reformidans, ad has malis recusabo persequeris, per etiam virtute invenire in.

Te nihil eruditi eam, sit aperiam accusam mediocritatem at. Nec ne nonumy dictas disputationi, vis ridens sadipscing ex. Harum euripidis ex vix, at consetetur instructior signiferumque mel, at mei elitr honestatis. Id sit congue vituperata. Temporibus eloquentiam no eum.

Pro id esse phaedrum, nostro iudicabit eos ut. Sit ea aperiam alienum, harum audiam voluptua cu usu. Iudico invenire te vel, id suscipit disputando pri. Ut sumo expetenda mea.

Cum at idque nullam aperiam, vis ex aeque ponderum luptatum. Vix soluta graeco dissentiet ut, ut est reque periculis similique, ut dicta dicant repudiare sea. Ne dolor legendos signiferumque ius, at eirmod convenire qui. Suas numquam conceptam mei ex. Autem homero eos et, sea dicta alienum iudicabit ut.

Ea duo consulatu vulputate, id elit perpetua cum. His ei aeque saepe audiam. Prompta laoreet facilisi ne sed, per hinc consetetur te, oratio fuisset ullamcorper mel at. Quis suscipiantur ne nec, agam efficiendi usu in.

Vis eu iuvaret singulis appellantur, usu ex saepe omittantur. Sed possit mnesarchum at, usu illum choro oratio in, et debet dolor vix. Mel aperiri suscipiantur ne, te per illum fuisset, lorem pericula mei ad. Pri id tale lucilius dissentiet, id sea sonet expetenda. Agam sensibus persequeris sed no, eum at tamquam sanctus.

Omnis exerci soleat ut vis. Rebum vidisse sea ex. Ius animal gubergren efficiantur ad, mollis probatus nec ut. Meis platonem ex vel, ut qui tale tritani equidem. Vide meis fuisset mel at, nam an assum delenit gubergren. No illum reprimique vim, te augue nullam per, ludus dicant suscipiantur ne sed.

An pri mediocrem deseruisse, ad sumo audire dissentiet sit. Sit ea civibus lobortis. Etiam ceteros commune ei vis. Pro ei equidem vivendo. Quo ne prima periculis omittantur, ex rebum veritus sit, ei dolor maiestatis mea.

\subsection{Lorem ipsum}

Et mel munere quodsi sapientem. Essent legimus ne pro. Est ornatus definiebas et. No habemus docendi ius, purto sapientem mei at. Tamquam vivendo necessitatibus has at, no habemus praesent nec. No quo modus iudicabit scriptorem. Modus intellegebat ea vim. Cu ius lorem regione offendit, ne accusata sensibus vituperatoribus quo. Sit ut iuvaret indoctum. Ut mea sale justo. Sapientem definitionem ius eu, at sea quem doming. Facete conclusionemque ut nec, vix at duis eius. Eos quot consequuntur et, ornatus liberavisse ne mei.

Per an dicam commodo tractatos, usu in timeam numquam tacimates. Case delectus eu sea, usu audiam eleifend tincidunt id, nec at decore discere mentitum. Ut elit veri eloquentiam his, ceteros tractatos ea has. Duo impetus scribentur et, eu quo errem everti, ad recusabo consulatu ius. Fastidii comprehensam pri ea, ex duo augue quando denique. Eos aeterno deserunt sententiae cu, ius quas tation patrioque ex.

Id autem scripta explicari nec, congue quidam possit te sit. Et usu ipsum bonorum graecis, ferri verear deterruisset eum cu. Purto porro accommodare cu vim. Cum ei tritani pertinacia voluptaria.



% itd.
% \appendix
% \include{dodatekA}
% \include{dodatekB}
% itd.

\printbibliography

\end{document}
